\chapter*{Abstract}
\addcontentsline{toc}{chapter}{Abstract}
\thispagestyle{plain}

Unemployment remains a major challenge in South Africa. In the second
quarter of 2026, the official unemployment rate was 33.6\%, with
approximately 8.5 million people unemployed \cite{statssa2026}.
Information and communication technologies (ICTs) cannot solve unemployment,
but they can be used to address some of the challenges experienced by job
seekers. This dissertation investigates the barriers experienced by job
seekers and explores how co-design can be used to develop an ICT intervention
to support them.

The research was conducted with job seekers and trainers at partner
organisations in Cape Town, South Africa. Interviews, workshops, prototyping
and a pilot were conducted through three iterative co-design cycles. The
cycles are described using Tell--Make--Enact. Participants first identified
challenges and developed initial ideas and paper prototypes. These were
reviewed and refined in the second cycle. The third cycle resulted in the
development of a mobile web-based artefact, which was piloted with
participants for two weeks. Exit interviews were conducted after the pilot.

The findings show that job seekers experienced several barriers, including
financial constraints, limited access to job information, skills and
language-related challenges, limited work experience, and exploitative or
discriminatory recruitment practices. Through the co-design process,
participants contributed to identifying these problems, suggesting solutions
and refining the intervention. The findings also show that while ICT can
address some barriers in the job-seeking process, it cannot address wider
structural problems such as the limited availability of jobs.

This dissertation shows how co-design can be used with job seekers in a
South African ICT4D context to understand their experiences and progressively
develop an intervention based on their needs and contributions. It also
highlights the importance of participation, context and intermediary support
when designing ICT interventions for job seekers.